\begin{remarks}[6.7.5]
\label{IV.2.s6.7.5-fr}
\begin{enumerate}
\item[(i)] En la demostración de que $\boldsymbol{P}(X,x)$ implica
$\boldsymbol{P}(X',x')$, no se puede prescindir de la hipótesis de que $k'$ sea una
extensión \emph{separable} de $k$. Ya se ha visto (4.6.1) cuando $\boldsymbol{P}$ es la
propiedad \emph{d)}; pero incluso cuando $X$ es \emph{geométricamente íntegro} (4.6.2),
puede ocurrir que $X$ sea \emph{regular} sin que $X'$ sea \emph{normal}.

Tomemos, por ejemplo, para $X$ una curva algebraica normal sobre $k$
(\textbf{II}, 7.4.2); como los anillos locales de $X$ son íntegramente cerrados y de
dimensión $1$ son anillos de valoración discreta, y por tanto regulares
(\textbf{II}, 7.1.6), y $X$ es, pues, un $k$-esquema regular (y \emph{con mayor razón}
verifica $(R_n)$ para todo $n\geqslant0$). Decir que $X$ es geométricamente
íntegro significa entonces que el cuerpo $K$ de las funciones racionales sobre
$X$ es una extensión separable y primaria de $k$ (4.6.3).

Ahora bien, tomemos para $k$ un cuerpo no perfecto de característica $p>2$, y
sea $a\in k$ un elemento que no pertenece a $k^p$. Sea
$B$ el anillo de polinomios $k[S,T]$ en dos indeterminadas $S$, $T$; el
polinomio
$P(S,T)=T^2\setmin S^p+a$
es irreducible en $B$, pues se verifica inmediatamente que $S^p\setmin a$ no es
un cuadrado en el anillo $k[S]$; por tanto, $X=\operatorname{Spec}(A)$, donde $A=B/PB$,
es una curva afín irreducible sobre $k$. Para ver que el esquema $X$
es regular, basta demostrar que es \emph{normal} (\textbf{II}, 7.4.5); ahora bien,
$A=k[S][t]$, donde $t$ es una raíz del polinomio $P$ considerado como
polinomio en $T$ sobre $k[S]$, de modo que el cuerpo $K=R(X)$ de las funciones
racionales sobre $X$ es el cuerpo de fracciones $k(S)[T]/(P)$ de $A$,
extensión cuadrática de $k(S)$, y por tanto \emph{separable} sobre $k(S)$, y \emph{con mayor
razón} sobre $k$. Como $2$ es invertible en $k$, se verifica inmediatamente que $A$
es el cierre integral de $k[S]$ en $K$, y es por tanto íntegramente cerrado,
lo que demuestra que $X$ es \emph{regular}. Por otra parte, si un elemento $f+tg$
de $K$ (con $f$, $g$ en $k(S)$) es algebraico sobre $k$, ocurre lo mismo con
su norma y su traza sobre $k(S)$, y como $k(S)$ es una extensión pura
de $k$, se concluye fácilmente que se debe tener $f\in k$ y $g=0$,
dicho de otro modo, $k$ es \emph{algebraicamente cerrado} en $K$, y \emph{con mayor razón} $K$
es una extensión \emph{primaria}
\oldpage[IV-2]{148}
de $k$ (4.3.1). Sin embargo, si
$k'=k(a^{1/p})$, $X'=X\otimes_k k'$ no es normal, pues en $k'[S]$
se puede escribir $S^p\setmin a=(S\setmin a^{1/p})^p$, y por tanto $X'$ es isomorfo a
$\operatorname{Spec}(A')$, donde $A'=k'[S][t']$, siendo $t'$ raíz del
polinomio $T^2\setmin S^p$. Ahora bien, $A'$ no es íntegramente cerrado, pues el elemento
$t''=t'/S$ del cuerpo de fracciones $K'=k'(S)[t']$ de $A'$ verifica
la ecuación de dependencia integral $t''^2\setmin S^{p-2}=0$ sobre $A'$ y
no pertenece a $A'$. En la teoría clásica, esto se expresa diciendo que el
«\emph{género}» sobre $k'$ del cuerpo de funciones racionales $K'$ de
$X'$ es \emph{estrictamente inferior} al de $K$ sobre $k$.

\item[(ii)] Como se ha recordado en (i), resulta de (4.6.1) que,
cuando $X$ es un preesquema algebraico sobre $k$ que no es
geométricamente reducido, existen extensiones radiciales finitas $k'$ de
$k$ tales que $X'=X\otimes_k k'$ no es reducido. Es interesante
dar un ejemplo de este hecho en el que $X$ es un esquema \emph{regular} sobre $k$, tal
que $k$ sea \emph{algebraicamente cerrado} en el cuerpo $K$ de las funciones
racionales de $X$. Sean $k$ un cuerpo de característica $p>0$ y $a$, $b$ que forman una familia $p$-libre sobre $k^p$. Designando de nuevo por $B$ el anillo $k[S,T]$, consideremos esta vez
el polinomio $P(S,T)=T^p\setmin aS^p\setmin b$; como $aS^p+b$ no es una potencia
$p$-ésima en $k(S)$, $P$ es irreducible como polinomio de $k(S)[T]$,
y el esquema $X_0=\operatorname{Spec}(B/PB)$ es, por tanto, una curva afín
\emph{íntegra} sobre $k$, cuyo cuerpo de funciones racionales es
$K=k(S)[t]$, donde $t$ es una raíz de $P$; demostremos que $k$ es
\emph{algebraicamente cerrado} en $K$. Supongamos, en efecto, que $K$ contiene un
elemento $z$ algebraico sobre $k$ y que no pertenece a $k$; se tendría entonces también
$z\notin k(S)$, y por tanto $[k(S)[z]:k(S)]>1$; como $[K:k(S)]=p$, se tendría $K=k(S)[z]$, y como $K$ es radicial sobre $k(S)$, se tendría
$z^p\in k(S)$, y por tanto $c=z^p\in k$, puesto que $z$ es algebraico sobre $k$;
pero se tendría $t^p=aS^p+b\in k^p(S^p)(c)$ y por tanto $a$ y $b$
pertenecerían a $k^p(c)$, lo cual es absurdo. Sea entonces $X$ la
\emph{normalización} de la curva $X_0$ en $K$, que es por tanto una curva normal
(y por consiguiente \emph{regular}) sobre $k$. Si $k'=k(a^{1/p},b^{1/p})$, está
claro que $aS^p+b$ es una potencia $p$-ésima en $k'(S)$, y por tanto
$K\otimes_k k'$ \emph{no es reducido}, ni \emph{con mayor razón} el esquema
$X\otimes_k k'$.
\end{enumerate}
\end{remarks}
